\section{前瞻：Testing-time Scaling、自改进与基础设施现实}

在展望部分，讲者给出 ``testing-time scaling'' 三个方向：让模型更深（deeper）、让执行更宽（wider / parallel thinking）、让执行更长（longer horizon, longer context, memory）。这与近年的推理时扩展路线高度一致。

\begin{importantbox}{讲者给出的未来判断}
未来高价值方向之一是 self-improvement：模型能够基于反馈持续改进自己。其难点不在单次优化，而在构建 ``弱模型生成更强模型'' 的稳定迭代机制。
\end{importantbox}

然而课程也非常坦诚地展示了现实困难：训练中的大量时间消耗在基础设施不稳定、作业失败恢复、资源调度与夜间值守上。讲者甚至用 ``90\% 时间不工作，10\% 时间真正推进'' 来形容训练体验。

\begin{warningbox}{被低估的最大成本：Infrastructure Instability}
GPU 掉线、通信失稳、作业重启、状态不可复现，会直接吞噬研究迭代速度。很多 ``算法慢'' 其实是 ``系统不稳''。
\end{warningbox}

\begin{knowledgebox}{把基础设施作为模型能力的一部分}
当训练链路规模达到千卡级别时，infra 工程质量会直接决定研究上限。稳定性、可观测性、容错恢复与自动化运维本质上属于 ``模型训练能力''。
\end{knowledgebox}

最后，讲者仍给出积极评价：尽管训练艰难，但当模型超过 baseline 并真实进入产品时，收益和成就感非常高。这也解释了为何工业界持续重投入 agentic model training。

\subsection{本章小结}
未来增益来自 ``更强 test-time scaling + 可持续自改进''，但落地前提是基础设施稳定。算法、系统、组织三者必须协同演进。

